% =====================================================================
% Main text for the beam search section.
% Companion to the algorithm listings alg:bs3d, alg:fch3d and the
% notation table tab:BS.
% =====================================================================

\subsection{Beam search}
\label{sec:bs}

Simulated annealing explores the space of \emph{loading sequences}: a move
perturbs the order in which items are offered, and the whole packing is then
rebuilt to discover what the perturbation achieved. The decoding of a sequence
into a packing therefore dominates the cost of an iteration, and the number of
iterations that fit into a given time budget is governed by it. This is a
property of the encoding rather than of the metaheuristic, so a second method
built on the same encoding would inherit the same limit.

Beam search changes the unit of work. Its states are \emph{partial packings}
rather than complete sequences, and a transition commits one further item to
one further space. The partial packing is extended incrementally and never
rebuilt, and the expensive construction that scores a state is amortised over
the entire subtree below it rather than over a single neighbour. We adopt the
beam search framework of \citet{xu2026} and adapt it to the MHKP; the
adaptations are described in Section~\ref{sec:bs-adapt}.

Beam search is a truncated breadth-first search. At every level the frontier is
scored by a heuristic evaluation function, and only the $w$ most promising
nodes are retained, the remainder being discarded. A second parameter $m$ caps
the number of children generated per node. Together, $w$ and $m$ bound the
width of the search tree and thus prevent the exponential growth that makes
exact tree search intractable, at the price of losing any optimality guarantee.
The notation is summarised in Table~\ref{tab:BS}.

\subsubsection{State representation and transition}

A state, or node, is a triple $N = (\mathcal{S}, P, \mathcal{B})$ consisting of
the free spaces $\mathcal{S}_j$ of every bin $j$, the items already placed, and
the items $\mathcal{B}$ still unplaced. In the root node $N_0$ each bin
contributes a single free space spanning its full extent
$(0,0,0,W_j,H_j,L_j)$, and no item is placed, so $\mathcal{B} = I$.

A transition is an \emph{action} $a = (j,s,i,o)$: item $i$ is placed at the
down--back--left (DBL) corner of free space $s$ of bin $j$ in rotation $o$.
Applying $a$ removes $i$ from $\mathcal{B}$ and re-partitions the free spaces
of bin $j$, as described in Section~\ref{sec:bs-spaces}. Only one position per
free space is considered, because the free-space representation already
enumerates the distinct positions at which an item can be seated. A state is
terminal when no bin has a free space that admits any remaining item.

Free spaces are maintained as \emph{maximal spaces} \citep{parreno2008}. When
an item is placed inside a space, that space is replaced by the residual slabs
on either side of the item along each of the three axes, each slab retaining
the full extent of the original space on the remaining two axes. At most six
slabs arise per space, and because each runs to a face of the original space
they overlap one another; any slab contained in another is discarded, which
keeps the list bounded. Retaining overlapping spaces rather than a disjoint
partition is deliberate: it guarantees that a later item can use the full
residual extent along any one axis, whereas a disjoint partition would divide
that extent arbitrarily among several smaller spaces and could reject an item
for which room physically remains. \textsc{UpdateFreeSpaces} in
Algorithm~\ref{alg:fch3d} states the procedure.
\label{sec:bs-spaces}

\subsubsection{Node expansion}

Generating every action available in a state would be prohibitive, since it
grows with the product of the number of free spaces, unplaced items and
rotations. \textsc{Expand} (Algorithm~\ref{alg:bs3d}) therefore restricts the
branching in three ways.

First, bins are traversed in non-increasing order of priority $p_j = 1/V_j$,
so that small bins are offered first. Because objective~(1) weights wasted
space by $p_j$, filling a small bin is worth as much as filling a large one,
and preferring small bins is what the objective rewards.

Second, within a bin the free spaces are traversed in DBL order, and the
procedure commits to the \emph{first} space that admits any item, returning
immediately. The search thus branches over which item to place next rather
than over where to place it. This is what keeps the branching factor
proportional to the number of items instead of to the number of
space--item--rotation triples, and the DBL order makes the choice a systematic
bottom-up filling of the bin rather than an arbitrary one.

Third, when the set $\mathcal{F}$ of items admissible in that space exceeds the
budget, a subset $\mathcal{F}'$ of $k = \lfloor m/|O| \rfloor$ items is drawn
by roulette selection with probabilities proportional to item volume,
$v_i / \sum_{i' \in \mathcal{F} \setminus \mathcal{F}'} v_{i'}$. Large items
are thus favoured, which reflects the fact that they are the harder ones to
accommodate later, while the randomisation retains a positive probability of
selecting any admissible item and so avoids the systematic blind spot of a
purely greedy rule. Crossing $\mathcal{F}'$ with the feasible rotations yields
at most $m$ actions, which is why the sample size is divided by $|O|$.

Membership in $\mathcal{F}$ requires more than geometric containment: an item
qualifies only if, in some rotation, it lies inside the space, overlaps no
placed item, and is stably supported in the sense of constraints
(10a)--(10r). Feasibility is thus enforced at expansion time, so that no
infeasible node ever enters the beam and no repair mechanism is required.

\subsubsection{Evaluation function}

A partial packing cannot be compared with another directly, since the volume
placed so far rewards nodes that are merely deeper. Each candidate node is
therefore \emph{completed} by the fast constructive heuristic \textsc{FCH} of
Algorithm~\ref{alg:fch3d}, and scored by the objective value of the resulting
complete solution. The score $g(N')$ thus estimates the quality of the best
solution reachable below $N'$, which is the quantity the beam should rank by.

\textsc{FCH} operates on a copy of the node, so that scoring never disturbs the
state it evaluates. It considers bins in priority order and, within a bin,
repeatedly performs the first feasible placement encountered when scanning free
spaces in DBL order and items in non-increasing volume. A bin is closed once no
placement remains, and the heuristic terminates when every bin is closed. The
resulting solution is scored by
\begin{equation}
f \;=\; \frac{\sum_{j \in J} p_j\, v_j(N')}{\sum_{j \in J} p_j\, V_j},
\label{eq:bs-score}
\end{equation}
the priority-weighted fill ratio. Expression~\eqref{eq:bs-score} is a strictly
decreasing affine transformation of objective~(1) normalised to $[0,1]$: it
induces exactly the same ranking of solutions and the same optimum, and merely
orients the objective so that larger values are better, as the beam requires.

Because \textsc{FCH} completes every candidate, it is executed once per
generated node and dominates the running time of the method. Its greedy,
first-fit design is therefore deliberate: the rollout must be fast enough that
the beam can afford one per candidate.

\subsubsection{The search}

\textsc{BeamSearch} (Algorithm~\ref{alg:bs3d}) assembles these components. The
beam is initialised with the root node, and \textsc{FCH}$(N_0)$ supplies an
incumbent solution $(P^\star, f^\star)$, so that a feasible solution is
available from the outset and at any moment thereafter. At each level, every
node of the beam $\mathcal{W}$ is expanded, each resulting action is applied to
obtain a child $N'$, and $N'$ is scored by a rollout. Since each rollout yields
a complete solution, the incumbent is updated whenever a rollout improves on
it; the best solution returned need not therefore correspond to any node
retained in the beam. All children of the level form the candidate set
$\mathcal{C}$, of which the $w$ highest-scoring are retained as the next beam.
The search terminates when the beam empties, that is when no node admits any
further placement, or when the time limit $T^{\max}$ is reached, and returns
the incumbent.

\subsubsection{Adaptation to the MHKP}
\label{sec:bs-adapt}

The framework of \citet{xu2026} addresses a single-bin problem with deformable
boxes, and four modifications are required here.

\emph{Multiple bins with priorities.} Their objective is the volume fraction of
one bin. Objective~(1) is the priority-weighted wasted space over several bins,
so a state carries the free spaces of every bin, an action names the bin it
applies to, and bins are offered in priority order.

\emph{Rotations.} Submodel V0 permits rotation about the vertical axis only,
and its items are rigid, so $|O| = 2$ rather than the eighteen
orientation--variant combinations of the original. The roulette budget is
$\lfloor m/2 \rfloor$ accordingly.

\emph{Stability.} Constraints (10a)--(10r) have no counterpart in the original
problem. They are verified at placement time, both during expansion and inside
\textsc{FCH}, so that every node of the search tree is feasible by
construction.

\emph{Scoring.} Nodes are ranked by the normalised complement
\eqref{eq:bs-score} rather than by objective~(1) directly, so that the beam's
selection of the $w$ largest values corresponds to the $w$ best solutions.

We retain the parameter values $w = 12$ and $m = 60$ obtained by
\citet{xu2026} through grid search. The two parameters govern the same
trade-offs here as there, width against depth and branching against rollout
cost, so their values provide a principled starting point; Section~\ref{sec:results}
reports a sensitivity analysis on our own instances.
